%!TEX root = ../main.tex

\begin{changemargin}{0.8cm}{0.8cm}

~\vfill{}

\section*{Abstrakt}
\vskip 0.5em

\todo{Check with someone czech :3}

Model-free \ac{RL} výrazně zlepšilo výkon autonomního závodění dronů, avšak stále trpí nízkou efektivitou využití vzorků.
Geometrické řízení poskytuje plynulost a agresivitu požadovanou při závodění, je však omezeno nutností ručního ladění regulačních zesílení.
Tato diplomová práce zkoumá, zda může přímé \ac{BPTT} zvýšit efektivitu využití vzorků oproti standardnímu algoritmu \ac{PPO} při zachování srovnatelného výkonu a zda se naučená adaptivní zesílení mohou ukázat jako účinnější než pevně nastavená zesílení v geometrickém řízení.
V prostředí PyTorch jsme vytvořili diferencovatelný simulátor kvadrokoptéry a porovnali různé trénovací úlohy i abstrakce řízení.
Zjistili jsme, že \ac{BPTT} je ve scénáři s jedinou bránou až $\mathbf{5\times}$ efektivnější z hlediska využití vzorků, avšak vykazuje chamtivé chování a na kompletních závodních tratích nepřekonává algoritmus \ac{PPO}.
Ani adaptivní zesílení v námi navržené formulaci nepřinesla lepší výsledky než zesílení pevně nastavená.
Přesto řadič s pevnými zesíleními optimalizovaný pomocí \ac{PPO} dosahuje výkonu srovnatelného se současným stavem poznání (state of the art), čímž náš simulátor představuje vhodný základ pro další výzkum.

\vskip 1em

{\bf Klíčová slova} \KlicovaSlova

\vskip 2.5cm

\end{changemargin}
